% ============================================================================
% 08-limitations.tex - Limitations (target 450--550 words)
% The four frozen research directions appear ONLY here (paper rule).
% Framing-only terms (VSM, rate-distortion) may appear here as explicit
% non-implementations; they must not appear as implemented mechanisms.
% ============================================================================

\section{Limitations}\label{sec:limitations}

\paragraph{Evidence scale, stated once.} Five limits bound every claim here, and
we state them once rather than scattering them. (i)~\emph{Parameters}: all have
provenance \texttt{guess}; the count of \texttt{measured} parameters is zero. The
phase-shift ratios ($0.70$/$0.85$), the trend threshold ($0.02$), the integral
gains, the ACT-R constants and the weight thresholds are design defaults, so no
threshold may be called tuned, validated or optimal. (ii)~\emph{Real traffic}:
eight real turns, five with a ratio observation, all at ratio $\approx 0$; the
calibration protocol's intake---thirty ratio-observed turns, ten boundary
samples, twenty result labels---is unmet. (iii)~\emph{Small samples}: the
external-baseline pilot is five rounds per cell on one host, one model route and
one day; the host pilot is ten cells, one run each; the mechanism supplement uses
twelve probes per class; the patch case study uses twenty payloads; each
held-out corpus has twenty items.
(iv)~\emph{Effect size}: apart from one pre-registered sign test on the ACT-R
class ($12/12$, $p=2.4\times10^{-4}$), no confidence intervals or $p$-values are
reported, and no forgetting, hallucination or end-to-end interception rate is
measured; proxies are labelled as proxies. The claim is reachability, not effect
size. (v)~\emph{Generalisation}: one host, one model route, tokens estimated from
characters; none of this extrapolates to another host, language or tokenizer.

\paragraph{External-baseline and patch-loop scope.} The three-subject comparison
(Section~\ref{sec:exp:baselines}) inherits the small-sample limit above, and its
subjects do not see the same input shape: the no-plugin host receives the raw
four-turn transcript, while the memory systems receive retrieved memories. The
Mem0 arm's out-of-the-box \texttt{add} path is additive, so the superseded value
is retained in the store unless a deletion is issued explicitly; the limits of
the comparison environment are stated separately below. The
interceptions we report for \system{} are \emph{structural} admission decisions:
the gate stops injection and records it, but the host still delivers the user
message, so the model may read the attack text. The pilot also exposed four
product findings that bound its readings: the fact extractor does not recognise
the natural correction phrasing (F-A) and stores the final question as a fact
(F-B); the tool gateway counts schema errors as blocks and injects hints into
later turns (F-C); and the first model request of a turn uses the previous
turn's injection snapshot, so a one-turn-stale memory can still reach the model
(F-D). The last one makes the S1 model-level answers non-interpretable and they
are marked n/i; the ledger-level supersession and a warm diagnostic are reported
separately. The red-team case study (Section~\ref{sec:exp:case-redteam}) is
in-sample by construction---the twenty payloads are the patch set---and the forty
held-out paraphrases stay at 0/40 because the patch is plugin-side while the
held-out evaluation is engine-side. The patch is a laboratory fork, not a product
change; its false-positive cost on benign traffic is not measured; and the
product architecture, including the absence of any adaptation module, is
unchanged.

\paragraph{Synthetic data and mock assertions.} The replay set is entirely
synthetic and permits a synthetic share of one, which disqualifies it for
deployment decisions. Plugin evidence is mock-host assertions
($188 + 148 + 45 + 15 = 396$) plus the ten-cell pilot above, not sustained
traffic. The trigger vocabularies are dominated by two languages, and token usage
is estimated from characters rather than produced by a tokenizer.

\paragraph{Chaos injection scope.} The chaos faults are synthetic single-shot
injections ($n=1$ per fault), not a sample of production failures. The
never-answering case is cut by an outer 150\,s window rather than a measured
client timeout; the warm state is pre-seeded; the context window is a laboratory
override (\texttt{LINGSHU\_CONTEXT\_WINDOW}) with anchors at section positions
rather than the absolute request head and tail; and one campaign used three
headless calls but fifteen provider requests because the agent loop expands a
call into several steps. Raw transcripts and per-field evidence files are kept,
and every claim in Section~\ref{sec:exp:chaos} is tied to a field rather than to a
summary sentence.

\paragraph{External-baseline environment limits.} The official \texttt{mem0ai}
2.2.1 library was run with a local non-neural hashed-bigram embedder (no
embedding API was available), and its default configuration produced empty
recall on cross-lingual queries. These are limits of our comparison environment,
not claims about Mem0.

\paragraph{Ablation power.} The pre-registered three-arm ablation has been
executed (zero-model, deterministic). Its original design supports only coarse
conclusions: the baseline--system retention difference is definitional, and
twelve probes detect only differences of about 0.5 or larger. A supplement
expands the ACT-R probes to twelve (Section~\ref{sec:exp:mech}), where the effect
becomes significant (12/12, $p=2.4\times10^{-4}$), while two mechanism arms
remain inseparable from the gate by construction. The gate itself points the
other way: zero of the forty held-out paraphrases (two corpora of twenty) are
intercepted (0/40, Section~\ref{sec:exp:heldout}), so lexical coverage is a
boundary, not a defence; the patch case study raises in-sample interception to
20/20 without moving this number.

\paragraph{Implementation debt.} Several mechanisms are implemented and unit
tested but not wired: impasse-driven subgoaling has no caller; the reduction path
has no conflict detector, so its trigger is always \texttt{none}; UCA-03 is
unreachable (refusal is never set); dynamic-weight eviction is never exercised
(the replay budget is never exceeded); the saturation-to-gate link is
caller-supplied. The assembled prompt contains no user text, so user-turn
jailbreaks cannot be caught engine-side; ACT-R implements base-level activation
only.

\paragraph{Interpretive framings are not implementations.} Two academic lenses
appear as \emph{framing-only}: the Viable System Model and a rate-distortion
analogy. Neither has a code locus, parameter or dependency: VSM does not occur in
the implementation and no rate--distortion functional is computed. They position
the work but cannot be cited as evidence.

\paragraph{Frozen research directions.} Four directions are explicitly
\emph{frozen}, \emph{not implemented}, and carry \emph{no engineering evidence}:
(i) action closure (closed-loop tool invocation with side-effect rollback);
(ii) adaptive evolution (meta-learning that rewrites its own search targets);
(iii) collective mind (a multi-agent blackboard); and (iv) multimodal
perception. No prototype, dependency or configuration switch exists
for any of them. Their absence bounds our claims: the control plane
governs prompt assembly, not tool execution, cross-agent coordination or
non-textual input.

\paragraph{Safety scope.} The fail-closed guarantee is behavioural and covers
injection only. We implement no independent crisis-safety pre-layer;
refusal depends on the host safety layer and a prompt-level constitution. We
therefore make no end-to-end safety claim: the structural comparisons of
Section~\ref{sec:exp:baselines} measure admission decisions (gate verdicts and
audit events) plus a small model-side pilot, not a safety rate that would
transfer to other hosts, models or attack distributions.
